\makesubsectionpage{Ensemble Learning}

\begin{frame}{Ensemble Learning: Grundidee}
\begin{columns}
\column{0.62\textwidth}
Einzelne Entscheidungsbäume neigen zur \textbf{Überanpassung}.

\vspace{0.4em}
\begin{itemize}
  \item \textbf{Idee:} mehrere verschiedene Bäume erstellen (ein
        \textbf{Wald}) und per \textbf{Mehrheitsentscheidung} klassifizieren.
  \item \textbf{Problem:} dafür braucht man \textbf{verschiedene} Bäume.
  \item \textbf{Lösung:} jeden Baum nur auf einer \textbf{zufälligen
        Teilmenge} der Trainingsdaten trainieren.
  \item[$\Rightarrow$] So entstehen lauter unterschiedliche Bäume.
\end{itemize}

\column{0.35\textwidth}
\centering
\twemoji[width=0.45\linewidth]{deciduous tree}\hfill
\twemoji[width=0.45\linewidth]{evergreen tree}
\end{columns}
\end{frame}

\begin{frame}{Ensemble Learning: Grundidee}
Ansätze, die mehrere Lerner kombinieren, nennt man \textbf{Ensemble
Learning}. Das Gesamtergebnis ist die (gewichtete) \textbf{Mehrheits\-%
entscheidung} der Basisklassifikatoren. Die zwei Varianten unterscheiden
sich darin, wie die Trainingsdaten für die einzelnen Lerner gebildet werden:

\vspace{0.6em}
\begin{columns}[T]
\column{0.48\textwidth}
\begin{blockC1}{Bagging}

Für jeden Lerner wird \textbf{zufällig} eine Teilmenge der Trainingsdaten
gezogen. Jeder Datensatz hat dieselbe Wahrscheinlichkeit.

\vspace{0.3em}
{\itshape Bootstrapping -- die Lerner entstehen unabhängig.}
\end{blockC1}

\column{0.48\textwidth}
\begin{blockC2}{Boosting}
Die Lerner entstehen \textbf{nacheinander}. Datensätze, die bisher
\textbf{falsch} vorhergesagt wurden, werden höher gewichtet.

\vspace{0.3em}
{\itshape Resampling/Reweighting -- die Lerner bauen aufeinander auf.}
\end{blockC2}
\end{columns}
\end{frame}

\begin{frame}{Ensemble Learning: Random Forests}
\begin{columns}
\column{0.45\textwidth}
Random Forests verwenden den \textbf{Bagging}-Ansatz. Für mehr Varianz
wird zusätzlich an jedem Entscheidungsknoten nur eine \textbf{zufällige
Teilmenge der Spalten} verwendet.

\vspace{0.4em}
Streicht man z.\,B. die Spalte mit der niedrigsten Entropie, rutscht eine
andere nach oben. Die einzelnen Bäume sind dadurch i.d.R. \textbf{schlechter}
als der vollständige Baum -- per \textbf{Mehrheitsentscheidung} wird das
Ergebnis aber meist \textbf{besser}.

\column{0.5\textwidth}
\centering
\resizebox{\linewidth}{!}{%
\begin{tikzpicture}[
  test/.style={draw=PrimaryColor, very thick, rounded corners=4pt,
               align=center, inner sep=4pt, text width=2.6cm,
               font=\scriptsize},
  ein/.style={fill=SecondaryColor, text=white, rounded corners=4pt,
              align=center, inner sep=4pt, text width=2.0cm,
              minimum height=0.9cm, font=\scriptsize},
  nein/.style={fill=ExampleTitle, text=white, rounded corners=4pt,
               align=center, inner sep=4pt, text width=2.0cm,
               minimum height=0.9cm, font=\scriptsize},
  kante/.style={->, >=Stealth, thick},
  lab/.style={font=\scriptsize\itshape, fill=white, inner sep=1pt}
]

% Wurzel (gestrichene Spalte)
\node[test] (r) at (0,0) {Bewerber*in hat Masterabschluss (90)};
\node[ein]  (e1) at (-2.8,-2.0) {Einstellen (40)};
\node[test] (b)  at (2.6,-2.0)  {sehr guter Bachelorabschluss (50)};
\node[ein]  (e2) at (0.4,-4.2)  {Einstellen (25)};
\node[test] (v)  at (4.4,-4.2)  {gut im Vorstellungsgespräch (25)};
\node[ein]  (e3) at (2.6,-6.4)  {Einstellen (20)};
\node[nein] (n)  at (6.2,-6.4)  {Nicht einstellen (20)};

\draw[kante] (r) -- (e1) node[lab, midway] {ja};
\draw[kante] (r) -- (b)  node[lab, midway] {nein};
\draw[kante] (b) -- (e2) node[lab, midway] {ja};
\draw[kante] (b) -- (v)  node[lab, midway] {nein};
\draw[kante] (v) -- (e3) node[lab, midway] {ja};
\draw[kante] (v) -- (n)  node[lab, midway] {nein};

% rotes X ueber die Wurzel (gestrichene Spalte)
\draw[red, line width=2pt] (r.south west) -- (r.north east);
\draw[red, line width=2pt] (r.north west) -- (r.south east);

\end{tikzpicture}}
\end{columns}
\end{frame}

\begin{frame}{Ensemble Learning: Boosting}
\begin{columns}
\column{0.60\textwidth}
Es gibt verschiedene \textbf{Boosting}-Verfahren. Sie unterscheiden sich
darin,
\begin{itemize}
  \item wie die Datenmengen für die schwachen Lerner gebildet werden und
  \item wie das \textbf{Gesamtergebnis} aus den einzelnen Lernern entsteht.
\end{itemize}

\vspace{0.4em}
Die gängigsten Verfahren sind \textbf{AdaBoost} und \textbf{LogitBoost}.

\column{0.33\textwidth}
\begin{blockC1}{In der Praxis}
Boosting mit \textbf{flachen} Entscheidungsbäumen als Basislerner gilt als
eines der \textbf{leistungsfähigsten} Klassifikationsmodelle -- oft genutzt
in der automatisierten \textbf{Bildanalyse}.
\end{blockC1}
\end{columns}
\source{\cite{guyon2011practical}}
\end{frame}

\makequizpage{\QUIZURL}{\QUIZQRCODE}{\QUIZIMAGE}
\makequestionspage{\FRAGENIMAGE}